\section{Related Work}

\subsection{Event-Centered EEG Analysis and Latency Variability}

EEG is often analyzed through events and latencies. Stimulus onset, response execution, error commission, movement onset, ERP component timing, and other temporally localized phenomena define not only what occurs in a trial, but also when task-relevant neural or behavioral processes unfold. This event-centered view is reflected in annotation practice: for example, Hierarchical Event Descriptors were introduced to describe laboratory and real-world EEG events in a machine-readable way for large-scale EEG analysis \citep{BigdelyShamlo2016HED}. It is also central to classical ERP analysis, where stimulus-locked and response-locked averaging are used to expose neural dynamics time-locked to task events. ERP-image visualization and related single-trial methods show that sorting trials by reaction time can reveal performance-linked temporal structure \citep{Jung1998}, while RIDE and subsequent reviews emphasize that intra-subject ERP latency variability can be estimated and can carry behaviorally meaningful information \citep{Ouyang2011RIDE,Ouyang2017}.

Response- and movement-related EEG provide additional examples of event-centered analysis. Lateralized readiness potentials have been used to study response preparation and response execution in time-resolved EEG \citep{Coles1989LRP}, and error-related negativities provide response-locked signatures of error commission and performance monitoring \citep{Gehring1993ErrorDetection,HolroydColes2002ERN}. Movement-related cortical potentials, including the Bereitschaftspotential, further show that EEG dynamics can be organized around movement onset and premovement preparation \citep{KornhuberDeecke1965BP,ShibasakiHallett2006BP}. This body of work motivates the broader premise that latency and event timing are not incidental details of EEG analysis; they are often part of the scientific object being studied.

This view is also compatible with decision-time accounts of behavior. Diffusion and accumulation-to-bound models explain RT variability through latent temporal dynamics in evidence accumulation and decision formation \citep{Ratcliff2008,OConnell2012,Kelly2013}. We do not fit a cognitive process model. Instead, we use the weaker premise that an observed RT is a noisy behavioral latency associated with an unobserved response-relevant event time.

Explicit EEG event detection is well established in settings where event annotations are available. Sleep-event detectors, for example, estimate when transient events such as spindles, K-complexes, or arousals occur rather than predicting only a window-level label \citep{Chambon2019DOSED}. These settings show that event detection is natural for EEG. However, they typically rely on explicit event annotations, such as expert-marked onsets, offsets, intervals, or external sensors. Our setting is different: we do not begin with manually annotated neural event times. We begin with a behavioral latency, RT, which is recorded as a scalar but implicitly provides a weak time-of-occurrence label. In this sense, our formulation is weakly supervised: the behavioral response provides timing evidence, but not a manually annotated neural event onset \citep{Zhou2018WeakSupervision}.

\subsection{Fixed-Window EEG Decoding of Latency-Structured Targets}

Modern supervised EEG decoding often maps a fixed stimulus-locked or task-locked EEG window to a scalar or class label. This formulation is convenient and aligns with standard benchmark metrics, but it can collapse temporal structure when the target itself has a time-of-occurrence interpretation. Reaction time is a clear example: it is evaluated as one scalar per trial, but it denotes the latency of a behavioral response following stimulus onset.

EEG-based RT estimation has often been benchmarked in this scalar-prediction form. Wu et al. used Riemannian tangent-space features for EEG-based BCI regression and evaluated RT estimation with RMSE and correlation \citep{Wu2017ReactionTimeRiemannian}. Chowdhury et al. later used periodogram features and deep neural networks to predict or classify visual stimulus-based RT from single-trial EEG, including a regression-based RT estimator \citep{Chowdhury2020DNNReactionTime}. The EEG Foundation Challenge provides a standardized scalar benchmark for trial-wise RT prediction from the HBN-EEG contrast change detection task, using normalized RMSE as the evaluation metric \citep{Aristimunha2025,EEGChallenge2025}. We keep this scalar evaluation protocol. Our question is complementary: because RT is a response latency inside the post-stimulus EEG window, can the same label be used internally as weak event-time supervision?

This distinction matters for interpretation. A scalar regressor may learn trial difficulty, subject-level response tendency, or stimulus-locked temporal priors without representing where the response-relevant event occurs inside the input window. A temporal readout alone does not fully resolve this issue: a model can predict weights over time bins and use a soft-argmax expectation while still being trained only through scalar RT error. Our experiments therefore separate direct scalar regression, temporal-readout regression, time-only segmentation, and distributionally supervised event-time modeling.

\subsection{Distributional and Time-to-Event Output Modeling}

From a machine-learning perspective, latency-structured EEG targets motivate output distributions rather than only point estimates. Label distribution learning replaces a single hard label with a distribution over related labels and is useful when neighboring labels share metric structure or when label ambiguity is meaningful \citep{Geng2016}. For RT, neighboring labels are ordered time bins rather than unrelated classes, making a smoothed event-time distribution a natural supervision target.

Time-to-event modeling provides a complementary probabilistic language. Classical survival analysis and modern neural time-to-event models estimate distributions over event times rather than only scalar predictions or risk scores \citep{Cox1972,Lee2018}. Distributional evaluation also makes calibration and uncertainty part of the modeling problem: posterior width, coverage, likelihood scores, and proper scoring rules can distinguish models with similar point errors but different uncertainty behavior \citep{Haider2020,Chapfuwa2023}.

We use this literature as output-distribution language, not as a claim that EEG RT prediction is a standard survival-analysis setting. We do not study long-horizon risk, censoring, or competing clinical outcomes. Each trial provides a finite post-stimulus EEG window and an observed behavioral latency. We therefore use the time-to-event view locally: the model predicts a posterior over response-relevant event times within the represented window, and the scalar RT prediction is obtained as the posterior expectation. EventNLL gives a probabilistic version of this formulation by treating observed RT as a noisy realization of a latent response-relevant event time.

\subsection{Relation to EEG Backbones, Pretraining, and Robust Decoding}

A separate line of EEG decoding work addresses cross-subject, cross-session, and cross-dataset variability through stronger input representations. Classical approaches include Riemannian covariance-based decoding and alignment methods for reducing subject or session mismatch \citep{Barachant2012,Jayaram2016,Zanini2018,HeWu2020}. Deep learning approaches pursue related goals through domain adaptation, moment matching, supervised EEG architectures, self-supervised pretraining, and foundation-style EEG backbones \citep{Ganin2016,Sun2016,Kostas2021,Wang2024,Jiang2024,Chen2024,Wang2025,Ouahidi2025,Yue2025}.

Our contribution is complementary to these input-representation approaches. We do not propose a new general-purpose EEG backbone, invariant representation loss, or pretraining strategy. Instead, we focus on the output representation and supervision. When the target is a latency, the label is not only a scalar value to regress; it is also a weak annotation of when a relevant event occurred. In our experiments, external EEG architectures and foundation-style models trained from scratch serve as controlled architecture baselines. The main question is whether aligning the learning objective with the event-time semantics of RT improves decoding and yields more interpretable trial-wise temporal evidence.

Taken together, prior work provides strong traditions for event-centered EEG analysis, fixed-window supervised decoding, explicit EEG event detection, and time-to-event modeling. To our knowledge, these traditions have not been systematically connected in EEG RT prediction through a controlled comparison between scalar regression, temporal-readout regression, and distributionally supervised event-time modeling. We fill this gap by treating behavioral latency as weak event-time supervision, learning a posterior over response-relevant latencies, and evaluating that posterior as temporal evidence rather than only as a source of a scalar RT estimate.
